\subsubsubsection{Block-DAGs Prevent DoS Attacks}

\label{sec:preventing-dos-attacks}

Consider the situation where an attacker is attempting to deny service via the production of empty blocks, and that the attacker can create blocks faster than the honest network.
Such a situation is illustrated in \autoref{fig:dag-dos-1}.
Since the goal of the attack is to prevent transactions from occurring, the attacker must\footnote{
    The attacker could also fill the blocks with spam transactions.
    That's more work for the attacker, but also more work for the honest network (calculating and storing that state, maybe indefinitely).
    It's preferable that the attacker has minimal transactions in their blocks.
    It's tempting to think of ideas like: \emph{since the attacker's blocks are empty, we can let honest nodes make larger blocks via some kind of weighted average block size calculation plus some flexibility in the size of blocks produced}.
    The problem with this is that it incents the attacker to fill their blocks with spam transactions, which is counterproductive.
} produce empty\footnote{
    Note that the attacker should still be reflecting other simplex-chains so as to maximize the total weight of the attacker's chain-segment.
    Given this, the attacker's blocks will contain reflected simplex-chain headers but no transactions.
} blocks.
Furthermore, if the attacker links back to any honest blocks then the honest blocks' histories will be merged with the canonical history; thus the attack fails.
The attacker's only available strategy is to produce a single chain of empty blocks.

\begin{figure}[ht]
    \centering
    \includegraphics[width=\linewidth]{dag_dos_v2_sag}
    \caption[
        An attempted empty-block DoS attack on a block-DAG.
    ]{
        An attempted empty-block DoS attack on a block-DAG.
        The attacker's blocks, $A_{i+\_}$, are mined in public and contain no transactions.
        Each block is annotated with its \emph{chain-weight} ($\Sigma_w$); a double outline indicates that a block is part of the main chain.
        Even though the attacker produces $2\times$ as many blocks as the honest network, the attack inevitably fails after a short while.
        Note: $H_i$ is defined to have $\Sigma_w = 0$ for illustrative convenience.
    }
    \label{fig:dag-dos-1}
\end{figure}

The challenge of such a DoS attack is to prevent honest miners from extending the attacker's chain-segment.
For traditional (non-DAG) chains --- where each non-genesis block has exactly one parent --- this is accomplished as soon as the attacker is able to reliably produce a heavier chain-segment than the honest network for a given period.\footnote{
    For a traditional blockchain, this would also imply a 51\% attack were possible, so this isn't a situation that those chains have to deal with.
    In a simplex, however, a miner can have $>50\%$ of a chain's local hash-rate, so we need to handle this case.
}

However, if blocks are allowed to have \emph{more} than one parent then there \emph{is no point} where an attacker can \emph{maintain} a DoS attack indefinitely.
Instead, they can only \emph{delay the execution} of some transactions for a short period of time.
Particularly, if an attacker can produce $A_{\text{blocks}} > 1$ for every 1 block produced by the honest network, then the attack can delay honest transactions for up to approximately $({A_{\text{blocks}}}^2 - 1) \cdot {B_f}^{-1}$ seconds, where $B_f$ is the frequency of block production (in Hz).
After this (approximate) point, the weight of the honest chain-segment, which includes the attacker's chain-segment, is always greatest.

If an attacker performs a \emph{repeating cycle} of these attacks, then they may be able to decrease the effective capacity of the chain by a factor of $\sim {A_{\text{blocks}}}^2$. %<!-- / (1 + A_{\text{blocks}}) -->
The opportunity cost of this attack, for the attacker, is at least as much as the lost transaction fees.
